\section{An exact series and two-sided bounds}
\label{a:sec:Abounds}

The product converts into a telescoping sum, and the sum is much more informative.
Write $r=2p$ and note that the recursion factorises as
\begin{equation*}
S_{n+1} = rS_n\lrb{1-\frac{S_n}{2}} .
\end{equation*}
Define $v_n = r^n/S_n$, so that $v_n\to1/A(p)$ by \cref{a:eq:Apdef}. Then
\begin{equation*}
v_{n+1}-v_n
= \frac{r^{n+1}}{rS_n\lrb{1-\frac{S_n}{2}}} - \frac{r^n}{S_n}
= \frac{r^n}{S_n}\lrbs{\frac{1}{1-\frac{S_n}{2}}-1}
= \frac{r^n}{S_n}\cdot\frac{S_n/2}{1-\frac{S_n}{2}}
= \frac{r^n}{2-S_n} .
\end{equation*}
Summing from $v_0 = 1/S_0 = 1$ gives an exact representation of the
quasi-stationary mean.

\begin{proposition}[Series representation]
\label{a:prop:series}
For $p\in[0,\tfrac12)$, with $S_n$ generated by \cref{a:eq:SnIter},
\begin{equation}
\label{a:eq:Aseries}
\frac{1}{A(p)} = 1 + \sum_{n=0}^{\infty}\frac{(2p)^n}{2-S_n} .
\end{equation}
\end{proposition}

\noindent
Unlike the product, this representation isolates the geometric decay in a factor
that does not depend on the dynamics at all; the whole of the nonlinearity sits in
the bounded denominator $2-S_n$. That is what makes it useful, because $S_n$ is
confined to $(0,1]$ and therefore $2-S_n$ to $[1,2)$. Bounding the denominator by
its extremes and summing the resulting geometric series gives the following.

\begin{proposition}[Two-sided bounds]
\label{a:prop:bounds}
With $\varepsilon=1-2p$ as in $\varepsilon=1-2p$, for $p\in(0,\tfrac12)$,
\begin{equation}
\label{a:eq:Abounds}
\frac{\varepsilon}{1+\varepsilon} \;<\; A(p) \;<\; \frac{2\varepsilon}{1+2\varepsilon},
\end{equation}
that is,
\begin{equation}
\label{a:eq:Aboundsexplicit}
\frac{1-2p}{2(1-p)} \;<\; A(p) \;<\; \frac{2(1-2p)}{3-4p} .
\end{equation}
The lower bound is exactly half the continuous-time constant \cref{a:eq:Acts}:
\begin{equation}
\label{a:eq:halfAc}
A(p) > \tfrac12\Ac(p) .
\end{equation}
\end{proposition}

\begin{proof}
Since $0<S_n\le1$ for all $n$ we have $1\le2-S_n<2$, and hence
$\tfrac12r^n < r^n/(2-S_n)\le r^n$, with equality on the right only at $n=0$, where
$S_0=1$. Summing over $n\ge0$ and using $\sum_nr^n = 1/(1-r) = 1/\varepsilon$,
\begin{equation*}
1+\frac{1}{2\varepsilon} \;<\; \frac{1}{A(p)} \;<\; 1+\frac{1}{\varepsilon},
\end{equation*}
and inverting gives \cref{a:eq:Abounds}. Substituting $\varepsilon=1-2p$ gives
\cref{a:eq:Aboundsexplicit}, whose lower bound is
$(1-2p)/\bigl(2(1-p)\bigr) = \tfrac12\Ac(p)$ by \cref{a:eq:Acts}.
\end{proof}

Two features of \cref{a:prop:bounds} deserve comment, because between them they
account for the whole shape of the discrete curve.

At the far subcritical end, $p\to0$, we have $\varepsilon\to1$ and the lower bound
tends to $\tfrac12$. Since $A(p)\le\tfrac12$ by \cref{a:prop:parity} below, the bound
is asymptotically attained and $A(0^+)=\tfrac12$. At the critical end,
$p\to\tfrac12^-$, we have $\varepsilon\to0$, the lower bound behaves like
$\varepsilon$ and the upper like $2\varepsilon$; \cref{a:sec:asymptotic} shows that
it is the upper bound that is attained there. So $A(p)$ runs from its lower bound
at one end of the interval to its upper bound at the other, which is why no
constant rescaling of $\Ac$ can reproduce it.

\begin{proposition}[Parity bound]
\label{a:prop:parity}
For the binary Galton--Watson process started from a single individual,
$A(p)\le\tfrac12$ for all $p\in[0,\tfrac12)$, with equality approached as $p\to0$.
\end{proposition}

\begin{proof}
Every individual leaves either $0$ or $2$ offspring, so $Z_n$ is even for every
$n\ge1$. Conditional on survival, therefore, $Z_n\ge2$, and hence
$\ex{Z_n\mid Z_n>0}\ge2$ for all $n\ge1$. Letting $n\to\infty$ and using
\cref{a:eq:qsmean} gives $1/A(p)\ge2$. As $p\to0$ the conditional population is $2$
with probability tending to one, so the bound is attained in the limit.
\end{proof}

\begin{table}[ht]
\centering
\caption{The constant $A(p)$ computed from the product \cref{a:eq:prodFormula} in
$60$-digit arithmetic, against the bounds of \cref{a:prop:bounds} and the
near-critical asymptotic $2(1-2p)$ of \cref{a:sec:asymptotic}. The bounds are loose
in the middle of the range and tighten at both ends: at $p=0.05$ the lower bound is
already within $3\%$ of $A$, whereas at $p=0.4999$ it is the upper bound that is
within $0.1\%$. The last two columns give the quasi-stationary mean and the number
of factors of \cref{a:eq:prodFormula} needed for the product to reach full precision
--- the cost that \cref{a:sec:practical} has to manage.}
\label{a:tab:Avals}
\begin{tabular}{ccccccr}
\toprule
$p$ & $\tfrac12\Ac(p)$ & $A(p)$ & $\dfrac{2(1-2p)}{3-4p}$ & $2(1-2p)$ & $1/A(p)$ & iterations\\
\midrule
0.05   & 0.473684 & 0.486180 & 0.642857 & 1.800000 & 2.0568 & 45 \\
0.1    & 0.444444 & 0.469382 & 0.615385 & 1.600000 & 2.1305 & 64 \\
0.2    & 0.375000 & 0.423895 & 0.545455 & 1.200000 & 2.3591 & 113 \\
0.3    & 0.285714 & 0.353967 & 0.444444 & 0.800000 & 2.8251 & 201 \\
0.4    & 0.166667 & 0.237647 & 0.285714 & 0.400000 & 4.2079 & 458 \\
0.45   & 0.090909 & 0.145069 & 0.166667 & 0.200000 & 6.8933 & 966 \\
0.48   & 0.038462 & 0.068011 & 0.074074 & 0.080000 & 14.7036 & 2\,473 \\
0.49   & 0.019608 & 0.036395 & 0.038462 & 0.040000 & 27.4764 & 4\,965 \\
0.499  & 0.001996 & 0.003944 & 0.003984 & 0.004000 & 253.519 & 48\,992 \\
0.4999 & 0.000200 & 0.000399 & 0.000400 & 0.000400 & 2504.65 & 478\,905 \\
\bottomrule
\end{tabular}
\end{table}

